\documentclass[UTF8,a4paper,12pt]{ctexbook}

\usepackage{amsmath,amssymb,bm}
\usepackage{geometry}
\usepackage{hyperref}

\geometry{margin=2.5cm}

\title{第2章\quad 向量自回归模型与统计推断基础}
\author{}
\date{}

\begin{document}

\maketitle

\setcounter{chapter}{1}
\chapter{向量自回归模型与统计推断基础}

\section{向量自回归（VAR）模型基础}

在多变量时间序列分析的量化实证研究中，正确且严谨地设定底层数据生成过程，是后续进行结构性假设检验、参数估计以及统计推断的先决条件。宏观经济运行指标、金融市场资产收益率以及各类微观主体的经济行为轨迹，往往不是在真空环境中孤立演化的，而是通过错综复杂的正负反馈机制、时间滞后效应以及同期交叉影响相互交织。这种跨越时间维度的动态相依性，要求研究者必须跳出传统单变量时间序列建模的局限，转而采用能够系统性刻画多变量联动关系的联立方程框架。

作为一种不依赖于严苛先验经济理论约束、而主要依托数据自身信息进行建模的多元时间序列工具，向量自回归（Vector Autoregression, VAR）模型之所以在计量经济学中占据重要地位，正是因为它通过将系统内所有内生变量的历史滞后项统一纳入解释变量集合，构造出一个全面、灵活且具有清晰代数结构的动态线性系统。该模型既能够描绘单个变量对其自身历史轨迹的依赖关系，也能够识别不同变量之间跨期传导与动态反馈的作用机制。因此，在宏观经济波动分析、货币政策传导研究、金融风险溢出检验以及资产收益联动建模等场景中，VAR 模型都构成了极为核心的分析框架。

本节将围绕 VAR 模型的基础理论展开讨论，重点介绍一般 VAR($p$) 模型的数学设定、参数矩阵的统计学含义、滞后算子的紧凑表示方法，以及模型平稳性的判定条件与特征方程机理。需要说明的是，以下内容主要聚焦于向量自回归模型本身的理论结构，而不涉及仿真实现、数值实验配置或具体算法部署等实现层面的内容。

\subsection{VAR($p$) 模型的数学设定}

向量自回归模型在本质上是经典单变量自回归过程（Autoregressive Process, AR）在多维欧几里得空间中的自然推广。它将单个标量随机变量的时间演化方程推广为一个包含多个相互作用、地位对等的内生变量所构成的向量差分系统。设一个包含 $N$ 个观测维度的多变量时间序列系统在任意离散时间点 $t$ 上的状态向量定义为
\begin{equation}
Y_t=(y_{1t},y_{2t},\dots,y_{Nt})^\top,
\end{equation}
其中，$Y_t$ 为一个 $N\times 1$ 的列向量，集中堆叠了系统在时刻 $t$ 的全部内生变量观测值。

若该系统当前状态不仅受到当期随机冲击的影响，而且还直接依赖于此前连续 $p$ 个时期的历史状态，则可将其表示为一个 $p$ 阶向量自回归模型，即 VAR($p$) 模型。其最一般且最标准的数学形式写为
\begin{equation}
Y_t = c + \Phi_1 Y_{t-1} + \Phi_2 Y_{t-2} + \cdots + \Phi_p Y_{t-p} + \varepsilon_t.
\end{equation}

在这个基本方程中，$c=(c_1,c_2,\dots,c_N)^\top$ 是一个 $N\times 1$ 的常数项向量，描述系统在均值意义上的基准位置；$\Phi_i$（$i=1,2,\dots,p$）是一系列维度为 $N\times N$ 的自回归系数矩阵；$\varepsilon_t=(\varepsilon_{1t},\varepsilon_{2t},\dots,\varepsilon_{Nt})^\top$ 为随机扰动向量，也常被称为新息过程（innovations）。

从统计解释来看，常数项向量 $c$ 对应于系统在不存在额外趋势项时的长期平均偏置，或者说是动态系统围绕其波动的均衡中心。更为关键的是系数矩阵 $\Phi_i$。它们共同构成了 VAR 模型的动力学核心。对于任意给定的矩阵元素 $\phi_{i,jk}$ 而言，该参数衡量的是第 $k$ 个变量在滞后 $i$ 期的取值变化，对第 $j$ 个变量当前值所产生的线性边际影响。沿着这一解释可以看到，系数矩阵的对角线元素反映各变量自身的历史惯性与持续性，而非对角线元素则刻画变量之间的跨期溢出与动态传导关系。

为了更清晰地揭示这种联立结构，可以将 VAR($p$) 的向量方程展开为 $N$ 个相互耦合的标量回归方程。对系统中的任意一个变量 $y_{it}$（$i=1,2,\dots,N$），其数据生成过程可写为
\begin{equation}
 y_{it} = c_i + \sum_{k=1}^{N}\phi_{1,ik}y_{k,t-1} + \sum_{k=1}^{N}\phi_{2,ik}y_{k,t-2} + \cdots + \sum_{k=1}^{N}\phi_{p,ik}y_{k,t-p} + \varepsilon_{it}.
\end{equation}

这一展开形式表明，系统中的每一个变量都由系统内所有变量在过去若干时期的历史取值共同解释。也就是说，在 VAR 模型中，每个方程的右侧都包含同一组解释变量，只是对应的系数不同。这种“所有变量解释所有变量”的对称设计，使模型不必在建模之初就对因果方向作出强约束，而是让变量之间的动态关系由数据本身加以揭示。这也正是 VAR 模型具备高度经验适用性的根本原因之一。

为了保证后续参数估计与统计推断的理论成立，通常需要对误差项施加一组标准弱白噪声条件。具体而言，扰动向量过程 $\{\varepsilon_t\}$ 一般被假定满足
\begin{align}
\mathbb{E}(\varepsilon_t) &= \bm{0},\\
\mathbb{E}(\varepsilon_t\varepsilon_t^\top) &= \Sigma,\\
\mathbb{E}(\varepsilon_t\varepsilon_s^\top) &= \bm{0}, \qquad t\neq s,
\end{align}
其中，$\Sigma$ 为一个对称正定的 $N\times N$ 协方差矩阵。上述条件意味着：第一，误差项在均值上不携带系统性偏移；第二，同一时点的扰动允许存在横截面相关性，但其协方差结构不随时间改变；第三，不同时点的误差项之间不应存在序列相关。由此，$\Sigma$ 的主对角线元素反映各变量所承受的特质性随机波动强度，而非对角线元素则度量不同变量在同一时点所遭遇的共同冲击或同期相关性。

需要强调的是，VAR 模型的这种灵活性并非没有代价。对于一个含有 $N$ 个变量、滞后阶数为 $p$ 的无约束 VAR 系统，仅自回归系数部分就需要估计 $N^2p$ 个参数；再加上 $N$ 个常数项以及协方差矩阵中 $N(N+1)/2$ 个自由参数，整个模型的参数规模会随着变量维度的扩大而迅速膨胀。因此，VAR 模型在低维环境中往往便于估计且解释明确，而在高维情形下则容易面临自由度不足、估计不稳定以及推断精度下降等问题。这一特征也解释了为什么后续文献中会进一步引入稀疏约束、低秩结构与正则化估计等高维扩展方法。不过，从模型基础出发，理解标准 VAR 的参数结构仍然是把握一切扩展形式的起点。

为了在理论推导中得到更紧凑的表达，时间序列文献通常引入滞后算子（lag operator）$L$。定义该算子满足
\begin{equation}
LY_t=Y_{t-1}, \qquad L^kY_t=Y_{t-k}.
\end{equation}
利用滞后算子，可以将原本冗长的 VAR($p$) 模型改写为矩阵多项式形式。定义
\begin{equation}
\Phi(L)=I_N-\Phi_1L-\Phi_2L^2-\cdots-\Phi_pL^p,
\end{equation}
则模型可被紧凑地写成
\begin{equation}
\Phi(L)Y_t = c + \varepsilon_t.
\end{equation}

这一算子表示方式在形式上大大简化了模型表达，更重要的是，它为后续平稳性分析、逆算子展开、脉冲响应函数以及向量移动平均表示的推导提供了统一的代数平台。因此，从研究方法论上看，滞后算子并不仅仅是一种书写简化技巧，而是连接动态系统表示与统计推断分析的重要工具。

\subsection{模型的平稳性条件与特征方程}

在建立 VAR($p$) 模型之后，必须进一步回答一个根本问题：该动态系统是否平稳。对时间序列分析而言，平稳性是参数估计、极限理论与预测分析成立的基础。如果一个随机过程的均值、方差及协方差结构会随着时间不断漂移，那么利用有限样本估计模型参数就很难获得稳定可靠的统计性质。对于多变量时间序列模型，文献中通常关注的是宽平稳性（weak stationarity），也即二阶矩意义下的平稳性。

设随机向量过程 $\{Y_t\}$ 的无条件均值为 $\mu$，则宽平稳要求其满足两个基本条件。第一，均值向量必须不依赖于时间，即
\begin{equation}
\mathbb{E}(Y_t)=\mu, \qquad \forall t.
\end{equation}
第二，任意两个时点之间的协方差矩阵只能依赖于它们的间隔 $h$，而不能依赖于绝对时点本身，即
\begin{equation}
\mathrm{Cov}(Y_t,Y_{t-h})=\Gamma_h, \qquad \forall t.
\end{equation}
同时，由协方差矩阵的对称性质可知，必然有
\begin{equation}
\Gamma_h=\Gamma_{-h}^\top.
\end{equation}

这意味着，一个宽平稳的 VAR 过程应当围绕一个固定均值波动，并且其动态依赖结构只与时间间隔有关，而不随样本所处的绝对位置改变。这样的性质为样本矩估计、一致性论证以及极限分布推导提供了必要前提。

然而，VAR($p$) 模型并不会自动满足平稳条件。其是否平稳，取决于系数矩阵 $\Phi_1,\Phi_2,\dots,\Phi_p$ 的整体组合方式。刻画这一稳定性的标准工具，是建立模型对应的特征方程。具体做法是将滞后算子形式中的 $L$ 替换为复变量 $z\in\mathbb{C}$，从而得到模型的特征矩阵多项式
\begin{equation}
\Phi(z)=I_N-\Phi_1 z-\Phi_2 z^2-\cdots-\Phi_p z^p.
\end{equation}

进一步地，令该矩阵的行列式为零，便得到 VAR($p$) 模型的特征方程：
\begin{equation}
\det\bigl(I_N-\Phi_1 z-\Phi_2 z^2-\cdots-\Phi_p z^p\bigr)=0.
\end{equation}

该方程本质上是一个关于复变量 $z$ 的高阶代数方程。它的全部根决定了系统的动态稳定性。根据多元时间序列分析的基本结论，VAR($p$) 模型宽平稳的充要条件是：上述特征方程的所有根都严格位于复平面的单位圆之外，即
\begin{equation}
|z_i|>1, \qquad i=1,2,\dots,Np.
\end{equation}

这一条件具有非常明确的经济学与统计学含义。如果全部根都位于单位圆外部，那么系统对任何有限冲击的反应都会随着时间推移逐渐衰减，冲击不会在系统中永久积累，因而序列的均值与方差结构能够保持稳定。相反，如果存在某个根落在单位圆上，则系统含有单位根，冲击的影响将具有永久性；如果存在某个根落在单位圆内部，则系统会出现发散或爆炸性行为，从而破坏一切基于稳定二阶矩结构的推断基础。

虽然特征方程给出了严格的理论判据，但在实际应用中，直接对高阶矩阵多项式求根往往并不方便。为此，时间序列理论中广泛采用伴随矩阵（companion matrix）表示，将一个高阶 VAR($p$) 系统转化为更高维的一阶状态空间表示。定义增广状态向量
\begin{equation}
X_t=
\begin{bmatrix}
Y_t\\
Y_{t-1}\\
\vdots\\
Y_{t-p+1}
\end{bmatrix},
\end{equation}
则 $X_t$ 的维度为 $Np\times 1$。在此基础上，原始 VAR($p$) 模型可重写为
\begin{equation}
X_t = \mathbf{C} + \mathbf{F}X_{t-1} + U_t,
\end{equation}
其中，$\mathbf{C}$ 为扩展常数项向量，$U_t$ 为扩展误差项，而 $\mathbf{F}$ 为系统的伴随矩阵。其标准分块形式为
\begin{equation}
\mathbf{F}=
\begin{bmatrix}
\Phi_1 & \Phi_2 & \cdots & \Phi_{p-1} & \Phi_p\\
I_N & 0 & \cdots & 0 & 0\\
0 & I_N & \cdots & 0 & 0\\
\vdots & \vdots & \ddots & \vdots & \vdots\\
0 & 0 & \cdots & I_N & 0
\end{bmatrix}.
\end{equation}

伴随矩阵方法的重要价值在于，它将原本复杂的高阶滞后系统稳定性问题，转化为标准矩阵谱分析问题。若记 $\lambda_1,\lambda_2,\dots,\lambda_{Np}$ 为伴随矩阵 $\mathbf{F}$ 的全部特征值，则这些特征值与原始特征方程的根之间满足倒数关系，即
\begin{equation}
\lambda_i = \frac{1}{z_i}.
\end{equation}
因此，VAR($p$) 的平稳条件可以等价写为
\begin{equation}
|\lambda_i|<1, \qquad i=1,2,\dots,Np,
\end{equation}
或者进一步压缩为更常用的谱半径条件：
\begin{equation}
\rho(\mathbf{F})=\max_{1\leq i\leq Np}|\lambda_i|<1.
\end{equation}

这一谱半径判据在理论与应用中都极其重要。理论上，它建立了特征方程、矩阵谱分解与动态稳定性之间的直接联系；应用上，它为判断 VAR 模型是否满足平稳性提供了简洁明了的数值标准。只要伴随矩阵的全部特征值都被严格限制在单位圆内部，便可认定模型具备稳定衰减性质。

当上述平稳条件成立时，VAR($p$) 模型还可以进一步转化为无穷阶向量移动平均表示，即 VMA($\infty$) 形式。设系统的无条件均值为 $\mu$，则平稳 VAR 过程可写成
\begin{equation}
Y_t = \mu + \sum_{i=0}^{\infty}\Theta_i\varepsilon_{t-i},
\end{equation}
其中 $\Theta_0=I_N$，而一系列矩阵系数 $\{\Theta_i\}_{i\ge 1}$ 则由原始自回归系数矩阵唯一决定。这一表示的含义是：当前的系统状态可以看作全部历史随机冲击经过动态传导后的加权叠加结果，而随着滞后期数增加，这些冲击的边际影响会逐步衰减。

VMA($\infty$) 表示不仅具有重要的理论意义，也构成了脉冲响应分析的基础。通过研究某一变量所受到的冲击如何经由系数矩阵在未来各期传递到系统中的其他变量，研究者可以进一步考察政策冲击、风险传染与多变量反馈机制的动态路径。由此可见，从 VAR($p$) 的基本向量方程、到滞后算子表示、再到特征方程、伴随矩阵与 VMA($\infty$) 表示，实际上构成了一条层层递进、彼此紧密衔接的理论链条。这些内容共同构成了向量自回归模型的基础知识框架，也为后续的参数估计、结构稳定性检验及更复杂的统计推断方法奠定了扎实基础。

\end{document}
